\usepackage{iftex}

\ifTUTeX
\usepackage{xltxtra}

\usepackage{fontspec}
\usepackage{polyglossia}
\setmainlanguage[babelshorthands=true]{russian}
\setotherlanguage{english}

\setmainfont{Times New Roman}
\setromanfont{Times New Roman} 
\setsansfont{Arial} 
\setmonofont{Courier New} 

\newfontfamily{\cyrillicfont}{Times New Roman} 
\newfontfamily{\cyrillicfontrm}{Times New Roman}
\newfontfamily{\cyrillicfonttt}{Courier New}
\newfontfamily{\cyrillicfontsf}{Arial}

% \usepackage{unicode-math}
\usepackage{amssymb}
\usepackage{amsfonts}

\else
% For pdfTeX we must use old
% 8-bit font technologies
\usepackage[T2A]{fontenc}
\usepackage[english,russian]{babel}
\usepackage[utf8]{inputenc}
\usepackage[english, russian]{babel}

%\usepackage{newtxmath}
\usepackage{amsmath}
\usepackage{amsthm}
\usepackage{amssymb}
\usepackage{amsfonts}
%\usepackage{mathtext}
\fi

%Hyphenation rules
\usepackage{hyphenat}
\hyphenation{ма-те-ма-ти-ка вос-ста-нав-ли-вать ото-бра-жа-ют-ся сге-не-ри-ро-ван-ном сим-во-лы}
% \usepackage[english, russian]{babel}

\usepackage{graphicx}
\graphicspath{ {docs/images/} {../images/} {./images/} }

\usepackage{tikz}
\usepackage{pgfplots} % для построения графиков
\usepackage{tkz-euclide}
\pgfplotsset{compat=1.18} % версия для совместимости
\usepackage{empheq}
\usetikzlibrary{positioning, calc, intersections}

\usepackage{float}

% \usepackage{comment}

\usepackage[hidelinks]{hyperref}

\newtheorem{theorem}{Теорема}[section]
\newtheorem{corollary}{Corollary}[theorem]
\newtheorem{lemma}[theorem]{Lemma}

\addto\captionsrussian{%
  \renewcommand{\figurename}{Рис.}%
  \renewcommand{\tablename}{Табл.}%
}

\usepackage{verbatim}

\newcommand{\cg}{\textnormal{\textsl{g}}}

\usepackage[dvipsnames]{xcolor}

\usepackage[siunitx, RPvoltages, american]{circuitikzgit}
\ctikzset{
    amplifiers/fill=cyan,
    sources/fill=green!20,
    diodes/fill=white,
    resistors/fill=violet,
    grounds/scale=1.2,
}
